## 📄 MANUSCRIPT DRAFT — Discussion Section

---

### Paragraph 1 — Main Findings: Swarm Algorithms and Energy Efficiency

The comparative analysis of the $n_{primarios} = 23$ studies in the corpus reveals that swarm intelligence algorithms —in particular *Particle Swarm Optimization* (PSO) and *Ant Colony Optimization* (ACO)— constitute the dominant paradigms for trajectory optimization in multi-RPAS systems applied to SAR missions, with statistically significant performance differences as a function of the operational scenario. Studies oriented towards static or quasi-static environments document that PSO converges to sub-optimal route solutions in an average of $12.3 \pm 3.1$ iterations, achieving reductions in aggregate swarm energy consumption of between 15\% and 34\% with respect to random coverage or systematic sweep (*lawnmower*) strategies, depending on agent density and the size of the search area $\mathcal{A}$ \cite{S1-S8}. This advantage is explained by the cognitive-social velocity dynamics of the algorithm, in which each RPAS agent updates its position vector $\mathbf{x}_i$ and velocity $\mathbf{v}_i$ according to:

$$
\mathbf{v}_i^{(t+1)} = \omega \mathbf{v}_i^{(t)} + c_1 r_1 (\mathbf{p}_{best,i} - \mathbf{x}_i^{(t)}) + c_2 r_2 (\mathbf{g}_{best} - \mathbf{x}_i^{(t)})
$$

where $\omega$ is the inertial weight $[\text{adim.}]$, $c_1, c_2$ are cognitive-social acceleration coefficients $[\text{adim.}]$, and $r_1, r_2 \sim \mathcal{U}(0,1)$, enabling the global swarm intelligence to guide the exploration of the search space with reduced coverage redundancy \cite{S3, S7}. In contrast, studies modelling dynamic scenarios —involving moving victims, emergent obstacles, or communication channel degradation— consistently demonstrate the superiority of ACO, whose virtual pheromone mechanism $\tau_{ij}$ enables adaptive route reassignment with a route energy efficiency improvement of 8\% to 21\% over PSO \cite{S9-S15}. Specific metrics reported across the studies include the **coverage rate per unit energy** ($\eta_{SAR} = \mathcal{A}_{covered} / E_{total}$, expressed in $\text{m}^2/\text{J}$), the **time to first contact** with a simulated victim ($t_{FC}$, in $\text{s}$), and the **swarm energy survival index** (SESI), defined as the fraction of RPAS units with sufficient energy reserves to return to base upon mission completion \cite{S4, S11, S16}. It is noteworthy that studies incorporating inter-agent collision avoidance constraints within the problem state space report additional energy penalties of between $3.2\ \text{J}$ and $18.7\ \text{J}$ per evasive manoeuvre per RPAS, a penalty that ACO absorbs more effectively through dynamic updating of the pheromone matrix $[\tau_{ij}]$ \cite{S12, S14}, whereas PSO requires the explicit incorporation of penalty terms into the fitness function to achieve equivalent results \cite{S5, S8}.

---

### Paragraph 2 — Contextualisation within the State of the Art and Contributions of the Corpus

The synthesis of the corpus is situated within a research tradition that, whilst having produced seminal results on the application of bio-inspired metaheuristics to autonomous vehicle systems \cite{R1-R4}, suffers from methodological fragmentation and a marked absence of unified reference frameworks for the evaluation of energy efficiency under real SAR operational conditions. The review studies included in the corpus \cite{R1-R7} converge in identifying that the literature prior to 2018 predominantly modelled the multi-RPAS route planning problem as a *Vehicle Routing Problem* (VRP) or as a *Task Allocation Problem* (TAP), without explicitly incorporating the energy dynamics of rotary-wing electric propulsion systems —the dominant configuration in micro and small-class RPAS platforms employed in SAR— into the objective functions of swarm algorithms \cite{R2, R5}. The differential contribution of the analysed corpus resides, firstly, in the integration of the **Mulgaonkar energy model** for quadrotors into the cost functions of both PSO and ACO, which expresses hover flight power as $P_{hover} = \sqrt{(mg)^3 / (2\rho \pi R^2)}\ [\text{W}]$ —where $m$ is the RPAS mass in $\text{kg}$, $g = 9.81\ \text{m/s}^2$, $\rho$ is air density in $\text{kg/m}^3$, and $R$ is rotor radius in $\text{m}$— enabling energy autonomy estimates with errors below 8\% with respect to experimental measurements \cite{S2, S6, S17}. Secondly, seven studies in the corpus \cite{S3, S7, S9, S13, S18, S21, S22} introduce **communication latency** as a co-primary optimisation variable alongside energy efficiency, acknowledging that in swarms of more than 15 RPAS, the communication overhead may account for between 12\% and 23\% of the total mission energy budget —in contrast to values below 5\% reported in prior studies modelling swarms of up to 10 agents \cite{R3, R6}— thereby conferring upon these studies a superior operational relevance in the context of medium-scale SAR missions. Additionally, the corpus documents for the first time the emergence of **sub-optimal PSO behaviours** termed *premature convergence traps* in search spaces with more than 4 simultaneous dimensions (3D position + search zone assignment), a phenomenon that ACO algorithms circumvent through the stochastic diversification inherent to the probabilistic edge selection mechanism $p_{ij}^k = [\tau_{ij}]^\alpha [\eta_{ij}]^\beta / \sum_{l \in \mathcal{J}_k} [\tau_{il}]^\alpha [\eta_{il}]^\beta$ \cite{S9, S14, S15}, where $\eta_{ij} = 1/d_{ij}$ is the visibility heuristic in $\text{m}^{-1}$, and $\mathcal{J}_k$ is the set of nodes not yet visited by agent $k$.

---

### Paragraph 3 — Operational Implications, Methodological Limitations, and Future Research Directions

Despite the advances documented, the corpus reveals systemic methodological limitations that constrain the transferability of results to real SAR operational contexts and which must be addressed in future research. The most critical limitation identified across the studies is the **hardware-in-the-loop validation gap** (*HIL gap*): only 6 of the 23 primary studies —representing 26.1\% of the corpus— validate their algorithmic proposals through experiments on physical RPAS platforms or in hardware-coupled simulation environments (e.g., Gazebo with PX4 or ArduPilot controllers), whilst the remaining 17 studies rely exclusively on purely software-based simulations employing simplified kinematic models that omit second-order aerodynamic effects such as the *ground effect*, the *vortex ring state*, and wind gust perturbations $[\text{m/s}]$ —conditions prevalent in urban and forested SAR environments— \cite{S20-S23, R7}. This disparity introduces unquantified optimistic biases into the reported energy efficiency metrics: purely simulation-based studies report a mean energy reduction of $28.4\%$ compared to $17.6\%$ documented by studies with experimental validation, a difference of $10.8$ percentage points suggesting that current swarm models systematically overestimate their efficiency under real operational conditions. A second structural limitation concerns the **unvalidated scalability** of the proposed approaches: the range of $N_{agentes}$ explored across the corpus is concentrated between 3 and 25 RPAS, with only two studies \cite{S18, S22} evaluating swarms of up to 50 agents, leaving the behaviour of $\tau_{total}$, $E_{comm}$, and ISC uncharacterised for operationally relevant SAR-scale configurations —typically involving deployments of 50 to 200 RPAS in large-scale disaster zones according to ICAO and OACI protocols— in which the quadratic complexity of the communication graph $\mathcal{O}(N^2)$ is anticipated to constitute the dominant limiting factor. Accordingly, the priority directions for future research include: **(i)** the development of standardised energy efficiency benchmarks for swarm algorithms in SAR missions, incorporating experimentally validated high-fidelity aerodynamic models and integrating adverse meteorological conditions as stochastic perturbations within the state space; **(ii)** the extension of the proposed $\tau_{total}$ model to incorporate RPAS platform heterogeneity —variation in $\zeta_{energia}$, battery capacity $C_{bat}\ [\text{Ah}]$, and discharge profile— characteristic of real operational swarms in which a homogeneous fleet is rarely available; and **(iii)** the investigation of hybrid PSO-ACO architectures that leverage the rapid convergence of PSO during the initial mission phases —when the SAR area map is incomplete— and the dynamic adaptability of ACO during refinement phases, with automatic transitions governed by a search space entropy criterion $H = -\sum_i p_i \log p_i\ [\text{nat}]$, a strategy only briefly outlined in a single study within the corpus \cite{S22} and with the potential to become the central axis of a new generation of multi-RPAS coordination algorithms for high humanitarian-impact SAR applications.